\section{원천 하나를 빼면 어떻게 되나}

\claimx{바깥 웹 말뭉치에서 들여온 학습 행을 통째로 빼도 유보 성능이 안 내려간다 --- 안쪽 두 원천은 빼면 내려간다}{바깥 원천을 뺀 팔의 델타가 잡음 문턱을 넘어 음수가 되거나, 안쪽 원천을 뺀 팔의 델타가 양수가 되면 이 주장은 떨어진다.}

학습 예산을 고정하고 바깥 원천의 섞음 비율만 바꾼 격자를 세웠다. 정칙화 세기는 자리마다 고르지 않고 등록된 값으로 얼렸다 --- 앞 노트에서는 그 값이 섞음 비율과 함께 움직여 두 효과가 엉켰다. 교차 유보 겹은 다섯 개의 씨앗에서 다시 만들어 평균했다. 바깥 원천을 통째로 뺀 팔의 델타는 $0.040949$ 이고, 다른 정칙화 값에서는 $0.011773$ 다. 두 값 모두 양수다. 같은 자리에서 안쪽 두 원천을 뺀 팔은 각각 $-0.026338$ 와 $-0.013606$ 로 음수다. 다만 도메인 균등 가중으로 재면 바깥 원천을 뺀 팔의 부호가 정칙화 값에 따라 갈린다($0.01344$ 와 $-0.000709$) --- 그래서 이 노트는 \textbf{빼면 나아진다}가 아니라 \textbf{넣어도 안 달라진다}까지만 주장한다.

\section{왜 그런가 --- 라벨을 부수는 대조}

\claimx{바깥 원천 쪽 학습 라벨을 통째로 뒤섞어도 유보 성능이 거의 안 움직이는 반면 안쪽 원천 쪽 라벨만 뒤섞으면 크게 움직인다 --- 즉 모형은 소수의 안쪽 행 위에 서 있다}{작은 쪽을 부순 절대 델타가 큰 쪽보다 작거나, 전량을 부순 델타가 잡음 문턱을 못 넘으면 이 주장은 떨어진다.}

학습 라벨만 뒤섞고 유보는 한 줄도 만지지 않았다. 바깥 원천 쪽 1,710 행을 전부 뒤섞은 절대 델타는 $0.018906$ 인데, 안쪽 원천 쪽 90 행만 뒤섞은 절대 델타는 $0.113826$ 다. 행당 효과로 고치면 $114.3919$ 배 차이다. 크기를 맞춘 대조 --- 바깥 원천에서 같은 수의 행만 뒤섞은 팔 --- 의 절대 델타는 $0.000916$ 로 사실상 영이다. 학습 라벨 전량을 뒤섞은 팔은 절대 델타 $0.252914$ 로 잡음 문턱 둘을 넘는다 --- 자가 학습 라벨을 보긴 본다.

\section{앞 노트의 크기는 대부분 정칙화 선택이었다}

\claimx{앞 노트가 원천 혼합의 효과로 적은 크기의 대부분은 정칙화 세기를 교차검증으로 다시 고른 데서 왔다}{정칙화를 맞춘 뒤에도 같은 크기가 남으면 이 주장은 떨어진다.}

앞 노트는 바깥 원천을 뺀 수준에서 델타 $0.127683$ 를 봤다. 정칙화를 등록된 값으로 맞추면 그 델타는 $0.011773$ 와 $0.040949$ 다. 정칙화 아홉 값 가운데 바깥 원천을 뺀 쪽이 나은 칸은 7 / 9 이다. 다만 정칙화를 각 팔에서 따로 최적화하면 두 팔의 최고값 차이는 $-0.000201$ 로 사실상 영이다 --- 그래서 두 문장을 갈라 적는다.

\section{자를 한 도메인에서 뗄 수 있나}

\claimx{묶음 가중 자는 한 도메인이 값의 대부분을 정하므로 도메인 균등 가중으로 갈아 끼워야 한다 --- 균등 가중은 학습 라벨을 전부 부수면 제대로 내려가므로 자로서 살아 있다}{균등 가중 값이 라벨 전량 파괴에서 안 내려가면 이 주장은 떨어진다.}

유보에서 가장 큰 도메인의 몫은 0.545994 인데 묶음 값에 대는 기여는 86.803 퍼센트다. 앞 노트의 격자에서 두 자의 부호가 뒤집히는 수준은 12 / 32 이다. 학습 라벨 전량을 뒤섞은 팔에서 균등 가중 델타는 $-0.023854$ 로 음수이므로, 측정 전에 박은 결정 규칙 셋이 모두 참이다. 이 노트는 균등 가중을 정본으로 올린다. \textbf{다만 등록한 조건은 부호만 물었다} --- 크기까지 물으면 그 델타는 자기 표준오차 $0.091314$ 의 두 배를 못 넘는다. 균등 가중은 묶음 가중보다 잡음이 넓고, 그 자의 검정력을 세우는 일은 다음 노트로 넘긴다. \textbf{올린 첫 결과는 이 노트 자신의 부차적 주장을 깎는 것이다} --- 정본 자로 보면 바깥 원천을 뺀 팔의 델타가 정칙화 값에 따라 부호를 바꾼다.
